%% Too Many Hands on the Throttle
%% IEEE conference submission draft (B1)
%% STATUS: complete formalism and architecture; Section IX is an evaluation
%% PROTOCOL, not results. Section VI's incident catalogue is a method plus an
%% empty table pending source verification.
%% See NOTES.md for open items and the citation verification list.

\documentclass[conference]{IEEEtran}
\IEEEoverridecommandlockouts

\usepackage{amsmath}
\usepackage{amssymb}
\usepackage{amsthm}
\usepackage{booktabs}
\usepackage{graphicx}
\usepackage{url}

\newtheorem{proposition}{Proposition}
\newtheorem{definition}{Definition}
\newtheorem{corollary}{Corollary}

\newcommand{\sys}{\textsc{Concerto}}

\begin{document}

\title{Too Many Hands on the Throttle: Arbitrating\\Overlapping Reliability Controls in Service Fleets}

\author{
\IEEEauthorblockN{Abhishek Pandey}
\IEEEauthorblockA{Meta\\
pandeyabhi1987@gmail.com}
\and
\IEEEauthorblockN{Abhinav Pandey}
\IEEEauthorblockA{Cisco\\
pandey.aby@gmail.com}
}

\maketitle

\begin{abstract}
A production service fleet runs an autoscaler, a load shedder, a circuit
breaker, a concurrency limiter, a rate limiter, a timeout policy, a retry
budget, and a router. These are separate control loops, owned by separate
teams, configured through separate surfaces, and tuned in isolation. They
all observe correlated signals and they all act on one plant. We show that
this arrangement is not merely untidy but structurally unsound, and that the
resulting failures are invisible to every controller involved.

We identify three coupling channels and characterize each. \emph{Actuation
coupling} arises when two actuators move the same plant variable: their
gains superpose, and we show that controllers individually tuned inside the
delay margin compose to an unstable loop, with the stability condition
depending on a quantity no single controller's configuration surface
exposes. \emph{Observation coupling} arises when one controller's action
corrupts another's input: we show that a fast load-reducing controller
acting upstream of a slow provisioning controller's measurement point
creates a positive feedback loop, and that detuning either controller
reduces the loop gain without removing the loop. \emph{Topological coupling}
arises when an action redistributes load, and we give the subcriticality
condition under which ejection does not cascade.

We then present \sys{}, an arbitration layer with one principle per channel:
a single severity estimate with a bounded authority budget allocated rather
than summed across actuators, counterfactual reconstruction of each
controller's input signal, and topology-gated admission of redistributive
actions. We establish that the premise holds in deployed software through an
audit of the configuration surfaces of Kubernetes, Envoy, Istio, and
Linkerd, none of which exposes a coordination hook between these mechanisms.

We report the model, the architecture, the configuration audit method, and a
full evaluation protocol over an open microservice testbed and a fleet
simulator. The experimental campaign is in progress and no measurements are
claimed here.
\end{abstract}

\begin{IEEEkeywords}
reliability control, autoscaling, load shedding, circuit breaking, control
allocation, cascading failure, microservices
\end{IEEEkeywords}

\section{Introduction}

Ask an operator how their service protects itself under load and you will
get a list. There is an autoscaler that adds replicas when utilization
rises. There is a load shedder that rejects low-priority work when latency
climbs. There is outlier detection that ejects unhealthy endpoints. There is
a concurrency limiter that caps in-flight requests, a rate limiter that caps
outbound calls, a timeout policy, a retry budget, and a router that moves
traffic away from trouble.

Each of these is a feedback controller. Each was designed in isolation,
documented in isolation, and tuned in isolation. Most are configured through
different surfaces, and in a large organization most are owned by different
teams. And every one of them observes some function of the same underlying
congestion state and acts, directly or indirectly, on that same state.

Control engineers have a name for this situation. When several actuators
have authority over one plant, the assignment of corrective effort among
them is the control allocation problem, and it has been treated rigorously
for decades in aerospace, where an aircraft has more control surfaces than
axes of motion \cite{durham1993,johansen2013}. Service fleets have the same
structure and do not use the same discipline. We allocate nothing. We run
every controller at full authority and hope.

This paper argues that the hope is misplaced in a specific and predictable
way, and that the resulting failure modes are invisible from inside any
individual controller.

\subsection{A failure no controller can see}

Consider a fleet with an autoscaler and a load shedder, both configured with
defaults that a careful operator would consider reasonable. Load rises. The
shedder sees latency climb and begins rejecting requests, which is exactly
its job, and latency recovers. The autoscaler, meanwhile, reads request rate
at the service, sees it fall, concludes that demand has dropped, and scales
the fleet down. The shedder now sees healthy latency on a smaller fleet and
relaxes. Full offered load returns to reduced capacity. Latency spikes, the
shedder clamps down harder, and the cycle repeats with a larger amplitude.

Every controller in this story behaved exactly as specified. The shedder shed
under latency pressure. The autoscaler scaled to observed demand. Neither
configuration is wrong. The fault is in the composition, and neither
controller's telemetry contains enough information to attribute it, because
each one sees only its own loop closing correctly.

\subsection{Contributions}

\begin{enumerate}
\item \textbf{A taxonomy of three coupling channels} (Section
\ref{sec:taxonomy}), distinguished by what is shared: the actuated variable,
the measured signal, or the load distribution. The distinction matters
because each channel requires a different remedy and two of the three are
untreatable by retuning.

\item \textbf{A formal characterization of each channel} (Section
\ref{sec:model}). Proposition \ref{prop:superpose} shows that controllers
individually tuned inside the delay margin compose to instability once their
margin fractions sum past unity, and Corollary \ref{cor:local} shows the
governing quantity is not observable at any single controller. Proposition
\ref{prop:masking} shows that signal masking creates a positive feedback
loop, and Corollary \ref{cor:detune} shows detuning cannot remove it.
Proposition \ref{prop:cascade} gives the subcriticality threshold for
redistributive actions.

\item \textbf{Evidence that the premise holds in deployed software} (Section
\ref{sec:evidence}). We audit the configuration surfaces of the mechanisms
operators actually run and record, for each, the signal it reads, the plant
variable it actuates, and whether any coordination hook with its peers
exists. We pair this with the incident record and with the practitioner
literature that already names these interactions.

\item \textbf{\sys{}, an arbitration layer} (Section \ref{sec:design}), with
one design principle per coupling channel plus a time-scale separation rule,
and an argument that each principle addresses its channel structurally
rather than by tuning.

\item \textbf{An evaluation protocol} (Section \ref{sec:protocol}) over an
open microservice testbed and a fleet simulator, with per-principle
ablations and an explicit statement of which baseline represents current
practice.
\end{enumerate}

We are explicit about scope. This paper delivers the taxonomy, the model,
the audit method, the architecture, and the protocol. The experimental
campaign is in progress. Section \ref{sec:protocol} states hypotheses and
measurement plans and reports no results.

\section{Background}
\label{sec:background}

\subsection{The mechanisms}

We use the following terms for the controllers a modern fleet runs. An
\emph{autoscaler} adjusts replica count from a utilization or demand signal.
A \emph{load shedder} or admission controller rejects a fraction of offered
work, usually prioritized, when a congestion signal exceeds a threshold. A
\emph{circuit breaker} or outlier detector removes an endpoint from the
serving set when its error rate or latency exceeds a threshold. A
\emph{concurrency limiter} caps in-flight requests. A \emph{rate limiter}
caps request rate. A \emph{timeout policy} bounds per-call latency. A
\emph{retry policy} reissues failed calls, typically bounded by a
\emph{retry budget}. A \emph{router} distributes traffic across endpoints and
performs failover.

\subsection{Why the composition is not studied}

Each mechanism has a substantial literature of its own. Autoscaling has
been treated as a control problem \cite{lorido2014}. Admission control and
overload protection have a long lineage in both networking and web serving
\cite{welsh2003}. Adaptive concurrency limiting draws on congestion control
\cite{netflixlimits}. Circuit breaking is folklore formalized into libraries.
Each literature assumes its mechanism is the one closing the loop.

The composition is discussed mainly in practitioner writing. The Google SRE
book devotes a chapter to cascading failures and to the interaction between
overload, retries, and capacity \cite{sre2016}. The Amazon Builders' Library
treats timeouts, retries, and backoff as a system-level concern and warns
about retry amplification \cite{buildersretry}. This literature is correct
and widely read. What it does not provide is a model that says when
composition is safe, which is what an operator needs in order to configure
more than one mechanism at a time.

\section{Three Coupling Channels}
\label{sec:taxonomy}

We distinguish coupling channels by what the controllers share.

\subsection{Actuation coupling}

Two controllers exhibit actuation coupling when their outputs move the same
plant variable in the same direction.

A concurrency limiter and a rate limiter both bound offered load. A shedder
and a rate limiter both bound admitted work. An autoscaler and an outlier
detector both change effective serving capacity, in opposite directions. In
each case the plant receives the sum of the corrections while each
controller was tuned as though it were applying the whole correction alone.

The orchestrator specification studied in our companion work states this as
an invariant it calls \emph{no double-tightening}: an implementation that
cuts a rate limit and reduces replica count for the same endpoint under mild
error is judged to have two controllers fighting one disturbance. That an
industrial specification names the failure is evidence that practitioners
encounter it.

\subsection{Observation coupling}

Two controllers exhibit observation coupling when one controller's actuation
lies upstream of the other's measurement point, so that the second
controller's input is a function of the first controller's output.

The autoscaler and shedder interaction of Section I is the canonical case.
Shedding removes work before it reaches the point where request rate is
measured, so the provisioning controller reads a demand signal that the
protection controller is actively suppressing. The same structure appears
whenever a protective mechanism is deployed at the edge of a measurement
boundary, which is the normal place to deploy one.

Observation coupling is the most dangerous of the three because it is the
least visible. Actuation coupling at least produces a plant that responds
too strongly, which an operator may recognize as overshoot. Observation
coupling produces a controller that is behaving correctly on false data.

\subsection{Topological coupling}

Two controllers exhibit topological coupling when one controller's action
redistributes load and thereby changes the state the other controller
observes at a different endpoint.

Ejection and failover routing are both redistributive. Removing an endpoint
does not remove its traffic; it moves it. If survivors were already near
saturation, the redistribution pushes them past the threshold that triggered
the first ejection, and the process compounds. Retry amplification is a
temporal variant: a timeout reduction converts slow successes into failures,
failures into retries, and retries into additional offered load, so a
tightening action increases the very quantity it was meant to reduce.

\section{A Model of Overlapping Authority}
\label{sec:model}

We model the fleet around an operating point and treat each coupling channel
in turn. The models are deliberately minimal. Their purpose is to identify
the structural condition under which composition fails, not to predict
trajectories.

\subsection{Setup}

Let $x(t)$ denote a scalar congestion state at an endpoint, for instance
queue depth or latency deviation from setpoint, normalized so that $x=0$ is
the operating point. Let controllers $i = 1, \ldots, m$ observe $x$ and
apply corrections $u_i$. Each controller has an actuation delay $\tau_i > 0$
comprising its measurement window, its decision interval, and the physical
time for its action to take effect. For an autoscaler $\tau$ includes
container start time and is often tens of seconds; for a shedder it is
milliseconds.

\subsection{Actuation coupling breaks the delay margin}

Consider $m$ controllers with overlapping authority, each applying
proportional feedback $u_i = -g_i x(t - \tau)$ with common delay $\tau$ and
gain $g_i > 0$. Acting alone, controller $i$ closes the loop
\[
\dot{x}(t) = -g_i\, x(t - \tau) + d(t) ,
\]
which is asymptotically stable if and only if $g_i \tau < \pi/2$, the
classical delay margin for first-order feedback with pure delay. Define
controller $i$'s \emph{margin fraction}
\[
\rho_i \;=\; \frac{g_i \tau}{\pi/2} \;\in\; (0,1) .
\]
A controller tuned to consume half its delay margin has $\rho_i = 1/2$.

\begin{proposition}[Gain superposition]
\label{prop:superpose}
If the $m$ controllers act concurrently on the same plant variable, the
composed loop is
\[
\dot{x}(t) = -\Big(\textstyle\sum_{i=1}^{m} g_i\Big)\, x(t - \tau) + d(t) ,
\]
which is unstable whenever $\sum_{i=1}^{m} \rho_i \geq 1$. In particular,
any two controllers each tuned to consume half or more of the delay margin
compose to an unstable loop, despite each being stable in isolation.
\end{proposition}

\begin{IEEEproof}[Proof sketch]
Superposition of the corrections gives the stated composed loop, whose
stability condition is $\left(\sum_i g_i\right)\tau < \pi/2$. Dividing by
$\pi/2$ yields $\sum_i \rho_i < 1$. Negating gives the claim.
\end{IEEEproof}

\begin{corollary}[Local tuning cannot prevent this]
\label{cor:local}
The margin fraction $\rho_i$ is determined by controller $i$'s own
configuration and the plant. The stability condition depends on
$\sum_j \rho_j$. No configuration surface enumerated in Section
\ref{sec:evidence} exposes $\rho_j$ for $j \neq i$, nor any proxy for it.
The condition is therefore not evaluable, let alone enforceable, at any
individual controller.
\end{corollary}

Corollary \ref{cor:local} is the practical heart of the actuation-coupling
argument. The failure is not that operators tune badly. It is that the
quantity governing stability is not exposed anywhere an operator could read
it. Adding a second well-tuned controller to a well-tuned system is an
unsafe operation, and nothing in the tooling says so.

Unequal delays complicate the algebra without changing the conclusion. When
$\tau_i$ differ, the composed characteristic equation
$s + \sum_i g_i e^{-s\tau_i} = 0$ replaces the single-delay form, and the
stability region shrinks relative to the single-controller case for the same
reason. We use the common-delay case for exposition.

\subsection{Observation coupling induces positive feedback}

Now consider two controllers with different roles. A fast controller $F$
reduces offered load with output $u_F \geq 0$, and a slow controller $S$
provisions capacity from a measured demand signal. Crucially, $F$ acts
upstream of $S$'s measurement point, so $S$ observes
\[
\hat{\lambda} \;=\; \lambda - u_F
\]
rather than the true offered load $\lambda$. Let capacity be
$\mu = \mu_0 + u_S$ with $u_S = g_S(\hat{\lambda} - \mu)$, let
$u_F = g_F\, x$ with $g_F > 0$, and let congestion respond to capacity with
sensitivity $\kappa = -\partial x / \partial \mu > 0$.

\begin{proposition}[Signal masking]
\label{prop:masking}
The composed system contains a feedback loop through the path
$x \to u_F \to \hat{\lambda} \to u_S \to \mu \to x$ whose loop gain is
$+\, g_F\, g_S\, \kappa$. The loop is positive, and the operating point is
unstable when $g_F g_S \kappa > 1$.
\end{proposition}

\begin{IEEEproof}[Proof sketch]
Trace the signs. An increase in congestion increases $u_F$ by $g_F > 0$. An
increase in $u_F$ decreases $\hat{\lambda}$ by construction. A decrease in
$\hat{\lambda}$ decreases $u_S$ by $g_S > 0$, hence decreases $\mu$. A
decrease in $\mu$ increases $x$ by $\kappa > 0$. The product of the four
sensitivities is positive, so the loop reinforces rather than opposes the
initial perturbation, and the standard small-gain argument gives instability
once the product exceeds unity.
\end{IEEEproof}

\begin{corollary}[Detuning does not remove the loop]
\label{cor:detune}
Reducing $g_F$ or $g_S$ reduces the loop gain but leaves the loop positive.
Any configuration with both controllers active and $F$ upstream of $S$'s
measurement retains a positive feedback path. Restoring negative feedback
requires severing the dependence of $\hat{\lambda}$ on $u_F$, which is a
structural change to the signal path rather than a tuning change.
\end{corollary}

Corollary \ref{cor:detune} explains an experience many operators report:
tuning the autoscaler's stabilization window makes the oscillation slower
without making it go away. The window changes $g_S$, and $g_S$ is not the
problem.

\subsection{Topological coupling has a branching threshold}

Let an endpoint be ejected and its load redistributed over the surviving
members of its pool. Let $\Gamma$ denote the expected number of additional
endpoints driven past the ejection threshold by the redistribution from one
ejection, given the current load margin distribution across the pool.

\begin{proposition}[Cascade subcriticality]
\label{prop:cascade}
Treating ejections as a branching process with mean offspring $\Gamma$, the
ejection cascade terminates almost surely if and only if $\Gamma < 1$. When
$\Gamma > 1$ the cascade reaches the pool's live-endpoint floor with positive
probability.
\end{proposition}

The practical content is that $\Gamma$ is computable before acting, from the
topology and the current per-endpoint headroom, which is what makes
topological coupling the most tractable of the three channels. The orchestrator
specification's cascade-risk invariant is precisely a bound on a quantity of
this form.

\section{Does the Premise Hold in Practice?}
\label{sec:evidence}

The model predicts trouble under two conditions: overlapping authority, and
independent configuration with no shared signal. This section establishes
that both conditions hold in software that is widely deployed. We use three
kinds of public evidence and deliberately avoid proprietary incident data,
so that every claim here is checkable by a reader.

\subsection{Configuration surface audit}

Method: for each mechanism, read the upstream documentation and
configuration reference, and record four fields. The signal the mechanism
reads. The plant variable it actuates. The tuning parameters it exposes.
Whether any documented hook allows it to observe or defer to a peer
mechanism. Table \ref{tab:audit} gives the schema and the mechanisms in
scope. We audit Kubernetes horizontal pod autoscaling and cluster
autoscaling, Envoy outlier detection, circuit breakers, adaptive concurrency,
local and global rate limiting, and retry budgets, Istio destination-rule
outlier detection, and Linkerd retry budgets.

\begin{table*}[t]
\caption{Configuration audit schema. Mechanisms in scope; the coordination-hook column is the finding the model predicts will be uniformly empty.}
\label{tab:audit}
\centering
\small
\begin{tabular}{@{}llll@{}}
\toprule
Mechanism & Reads & Actuates & Coord. hook \\
\midrule
K8s HPA & utilization & replicas & --- \\
K8s cluster autoscaler & pending pods & nodes & --- \\
Envoy outlier detection & errors, latency & serving set & --- \\
Envoy circuit breakers & in-flight, pending & admitted load & --- \\
Envoy adaptive concurrency & latency & concurrency & --- \\
Envoy rate limiting & request rate & admitted load & --- \\
Envoy retry budget & failure rate & offered load & --- \\
Istio outlier detection & errors & serving set & --- \\
Linkerd retry budget & failure rate & offered load & --- \\
\bottomrule
\end{tabular}
\end{table*}

We state the expected finding as a hypothesis rather than a result, since
the audit is in progress: we expect the coordination-hook column to be
uniformly empty, and we expect at least three pairs of mechanisms in the
table to actuate the same plant variable. If either expectation fails the
paper's premise weakens, and we will report that.

\subsection{Incident catalogue}

Public postmortems establish that these interactions occur outside
simulation. We use them for existence and mechanism, not for prevalence.
This is a deliberate methodological limit. Vendors publish selectively, and
most writeups lack the detail needed to diagnose a controller interaction,
so any rate computed over published incidents would be an artifact of
disclosure policy rather than a property of systems.

Method: rather than inherit a third-party curated list, whose maintenance we
cannot depend on, we index primary sources directly. We search vendor
incident archives from major cloud and infrastructure providers, apply stated
inclusion criteria, and code each admitted incident against the three
channels of Section \ref{sec:taxonomy}. Both authors code independently and
we report agreement. The resulting index, with inclusion criteria and per
incident coding, is released as an artifact so that the catalogue does not
inherit our own maintenance either.

We also draw on research corpora that have publications behind them and are
therefore more durable than community lists, including published studies of
cloud fault life cycles and of user- and provider-reported cloud failures.

\subsection{Practitioner literature}

Independent of any dataset, the interactions we model are already named in
the canonical operational literature. The Google SRE treatment of cascading
failures describes overload propagation and the role of retries in
amplifying it \cite{sre2016}. The Amazon Builders' Library treats retry
amplification and load shedding as system-level concerns rather than local
ones \cite{buildersretry,buildersshed}. We cite this literature as evidence
that the phenomena are recognized, and position our contribution as supplying
the model and the remedy that it does not.

\section{\sys{}: Arbitration by Construction}
\label{sec:design}

\sys{} is an arbitration layer that sits between the controllers and the
plant. Controllers become advisors: they compute desired corrections, and
\sys{} decides what is actually applied. The name reflects the intent, which
is controllers playing from one score rather than each at full volume.

The design has one principle per coupling channel, plus a cross-cutting
scheduling rule.

\subsection{P1: bounded shared authority}

\emph{Addresses actuation coupling.}

\sys{} computes a single severity estimate $s \in [0,1]$ from the plant
state, once, and derives a total corrective authority budget $B(s)$. Each
actuator receives an allocation $\alpha_i(s)$ subject to
$\sum_i \alpha_i(s) \leq B(s)$. Allocations are assigned by actuation cost
and time constant, not by letting each controller take what it wants.

The structural property is that total loop gain is bounded by $B$ regardless
of how many controllers are enabled. Enabling an additional controller
redistributes authority rather than adding to it, so the margin-fraction sum
of Proposition \ref{prop:superpose} is bounded by construction and
Corollary \ref{cor:local} no longer bites: the governing quantity is
computed in one place, by the component that owns it.

This is control allocation in the aerospace sense \cite{durham1993}, applied
to a plant whose actuators happen to be named autoscaler and shedder.

\subsection{P2: counterfactual signal reconstruction}

\emph{Addresses observation coupling.}

Each controller is presented with the signal it would observe if no other
controller were acting. Concretely, the provisioning controller reads
$\lambda = \hat{\lambda} + u_F$, the sum of admitted and shed load, rather
than admitted load alone. Generalizing, \sys{} maintains for each exposed
signal an accounting of the corrections applied upstream of its measurement
point, and serves the reconstructed pre-intervention value.

By Corollary \ref{cor:detune} this is the only class of remedy that works,
because it removes the dependence of $\hat{\lambda}$ on $u_F$ rather than
shrinking it. The implementation cost is low: sheddders and rate limiters
already count rejected work for their own telemetry, so the required
quantity usually exists and is merely not wired to the consumer that needs
it.

\subsection{P3: topology-gated action admission}

\emph{Addresses topological coupling.}

Before applying any redistributive action, \sys{} computes the projected
cascade gain $\Gamma$ of Proposition \ref{prop:cascade} from the current
topology and per-endpoint headroom, and admits the action only if the
post-action $\Gamma$ remains below a subcriticality threshold. When the full
requested action set would exceed it, the set is scaled back rather than
dropped, preserving as much protective value as the threshold allows.

\subsection{P4: time-scale separation}

\emph{Cross-cutting.}

Even under a shared budget, two actuators can chase the same disturbance
component. \sys{} splits the severity signal by frequency band and assigns
bands to actuators by time constant: the fast band drives shedding and
concurrency, the slow band drives replica provisioning. An actuator never
receives authority over a disturbance component faster than its own
actuation delay, which prevents a slow controller from reacting to transients
it cannot affect in time, a known source of autoscaler thrash.

\section{Implementation}
\label{sec:impl}

We implement \sys{} in two settings.

The first is a sidecar and controller pair for Kubernetes and Envoy. The
severity estimator and allocator run as a cluster-scoped controller. Existing
mechanisms are placed in advisory mode where their configuration permits it,
and where it does not, they are bounded externally by the allocator setting
their limits each interval. Signal reconstruction is implemented by exposing
a reconstructed demand metric that the autoscaler consumes in place of the
raw one, which requires no change to the autoscaler itself.

The second is a reference implementation inside a fleet simulator, used for
controlled ablation at scale where testbed experiments are impractical. The
simulator exposes the same public decision surface used in our companion
work on grading long-horizon controllers, which lets the same arbitration
logic be evaluated under both.

A deliberate design constraint is that \sys{} requires no modification to the
controllers it arbitrates. Each of the three principles is implementable by
bounding a controller's output, substituting its input, or gating its action,
all of which are available from outside.

\section{Evaluation Protocol}
\label{sec:protocol}

This section states the campaign we are running. It reports no results.

\subsection{Testbed}

An open microservice benchmark on Kubernetes with Envoy or Istio providing
outlier detection, retries, and concurrency limiting, and horizontal pod
autoscaling providing capacity. We favor DeathStarBench as the primary
application for its established use in the systems community
\cite{gan2019}, with a lighter alternative retained for iteration speed.
Workloads are driven by an open-loop generator, with arrival patterns taken
from public cluster traces to obtain realistic non-stationarity.

The testbed carries the paper's empirical weight. The central experiment is
the observation-coupling demonstration, because it is the channel that
retuning cannot address: we aim to produce the autoscaler and shedder
oscillation of Section I on unmodified, reasonably configured, widely
deployed software.

\subsection{Simulator}

A fleet simulator supplies controlled ablation across many scenarios and at
scales the testbed cannot reach. It is the instrument for isolating
mechanism, not for establishing that the phenomenon is real; the testbed does
that.

\subsection{Baselines}

\textbf{B1, independent defaults.} Every mechanism enabled with vendor
default configuration. This is not a strawman. It is what a team gets by
following each mechanism's own documentation.

\textbf{B2, independently retuned.} Each controller detuned by an expert to
compensate for observed instability, without any shared signal. This
baseline tests whether the problem is a tuning problem, and Corollary
\ref{cor:detune} predicts it will fail specifically on observation coupling
while improving actuation coupling.

\textbf{B3, strict precedence arbiter.} A naive arbitration scheme in which a
fixed priority order determines which controller acts when several wish to.
This isolates the value of proportional allocation over simple mutual
exclusion.

\textbf{B4, \sys{}.}

\subsection{Metrics and research questions}

Metrics: oscillation amplitude and period after a step disturbance, settling
time, objective attainment, capacity cost, cascade incidence, goodput under
a retry storm, and the rate of simultaneous tightening actions on one
endpoint.

\textbf{RQ1.} Does each coupling channel reproduce on unmodified deployed
software under reasonable configuration?

\textbf{RQ2.} Does expert retuning remove instability from actuation coupling
but not from observation coupling, as Corollary \ref{cor:detune} predicts?
This is the sharpest falsifiable prediction the model makes and we treat it
as the paper's central test.

\textbf{RQ3.} Does \sys{} restore stability without a cost or attainment
penalty relative to the best-performing baseline?

\textbf{RQ4.} What does each principle contribute? Measured by removing P1
through P4 individually.

\textbf{RQ5.} How does arbitration behave as the number of enabled mechanisms
grows? Proposition \ref{prop:superpose} predicts the independent baselines
degrade with $m$ while \sys{} does not.

\begin{table}[t]
\caption{Results table structure. No values are reported; the campaign is in progress.}
\label{tab:placeholder}
\centering
\small
\begin{tabular}{@{}lcccc@{}}
\toprule
Configuration & Osc. & Settle & Attain. & Cost \\
\midrule
B1 defaults & --- & --- & --- & --- \\
B2 retuned & --- & --- & --- & --- \\
B3 precedence & --- & --- & --- & --- \\
B4 \sys{} & --- & --- & --- & --- \\
\;\;\textit{less} P1 & --- & --- & --- & --- \\
\;\;\textit{less} P2 & --- & --- & --- & --- \\
\;\;\textit{less} P3 & --- & --- & --- & --- \\
\;\;\textit{less} P4 & --- & --- & --- & --- \\
\bottomrule
\end{tabular}
\end{table}

\section{Related Work}
\label{sec:related}

\textbf{Control allocation.} The assignment of corrective effort among
redundant actuators is well developed in aerospace control
\cite{durham1993,johansen2013}. Our contribution is not to that theory but
to the observation that service fleets have the same actuator redundancy and
none of the allocation discipline, and to the two coupling channels,
observational and topological, that the aerospace formulation does not
address because aircraft control surfaces do not corrupt each other's
sensors or redistribute load across a graph.

\textbf{Autoscaling.} Autoscaling has been studied extensively as a control
and prediction problem \cite{lorido2014}. This work assumes the autoscaler
is the controller closing the loop. We study what happens when it is not the
only one, and Proposition \ref{prop:masking} concerns the validity of the
demand signal such work takes as given.

\textbf{Overload control and admission.} Overload protection for servers has
a long lineage \cite{welsh2003}, and adaptive concurrency limiting imports
congestion-control ideas into the application layer
\cite{netflixlimits}. Again the composition is out of scope in that
literature.

\textbf{Cascading failure.} Cascade dynamics have been modeled in power grids
and in networks, and the branching-process threshold of Proposition
\ref{prop:cascade} is standard there. The contribution here is to treat
cascade risk as an admission gate on a controller's action rather than as a
post-hoc analysis.

\textbf{Practitioner literature.} The operational canon
\cite{sre2016,buildersretry,buildersshed} names these interactions and
prescribes remedies per mechanism. We supply the composition model and an
arbitration layer, which it does not.

\section{Threats to Validity}
\label{sec:threats}

\textbf{Model fidelity.} The models of Section \ref{sec:model} are linearized
and scalar. They are intended to identify structural conditions, not to
predict trajectories, and the testbed is the check on whether the conditions
they identify actually govern real systems.

\textbf{Configuration realism.} A reviewer may object that a competent
operator would not deploy the configurations in baseline B1. The audit of
Section \ref{sec:evidence} is partly a response: we report what the defaults
are and what the documentation recommends, so that the baseline is defensible
by citation rather than by assertion.

\textbf{Benchmark representativeness.} Findings on a benchmark application
may not transfer to production topologies, which are larger and more
irregular. The simulator partially addresses scale and the audit addresses
mechanism coverage, but neither substitutes for production validation, which
we do not have.

\textbf{Incident catalogue bias.} Addressed directly in Section
\ref{sec:evidence}. We make no prevalence claim, and the catalogue supports
existence and mechanism only.

\section{Conclusion}

Service fleets run many reliability controllers over one plant and allocate
authority among them not at all. We identified three channels through which
those controllers interfere, and showed that each has a distinct structural
signature: summed gains that break a delay margin no controller can observe,
a masked signal that turns negative feedback positive and cannot be retuned
away, and a redistribution process with a branching threshold. We presented
an arbitration layer with one principle per channel, implementable without
modifying the controllers it governs.

The broader claim is that reliability mechanisms should be designed for
composition and are not. Every mechanism in the audit of Section
\ref{sec:evidence} is individually sound and individually well documented,
and not one of them exposes the information a neighboring controller would
need in order to be safe. That is a gap in how the mechanisms are specified,
not in how operators use them.



\begin{thebibliography}{99}
\bibitem{durham1993} W. C. Durham, ``Constrained control allocation,'' \emph{Journal of Guidance, Control, and Dynamics}, vol. 16, no. 4, pp. 717--725, 1993.
\bibitem{johansen2013} T. A. Johansen and T. I. Fossen, ``Control allocation: A survey,'' \emph{Automatica}, vol. 49, no. 5, pp. 1087--1103, 2013.
\bibitem{lorido2014} T. Lorido-Botr\'an, J. Miguel-Alonso, and J. A. Lozano, ``A review of auto-scaling techniques for elastic applications in cloud environments,'' \emph{Journal of Grid Computing}, vol. 12, no. 4, pp. 559--592, 2014.
\bibitem{welsh2003} M. Welsh and D. Culler, ``Adaptive overload control for busy internet servers,'' in \emph{USENIX Symposium on Internet Technologies and Systems (USITS)}, 2003.
\bibitem{netflixlimits} Netflix Technology Blog, ``Performance under load: Adaptive concurrency limits,'' 2018. [Online]. Available: \url{https://netflixtechblog.medium.com/}
\bibitem{sre2016} B. Beyer, C. Jones, J. Petoff, and N. R. Murphy, Eds., \emph{Site Reliability Engineering: How Google Runs Production Systems}. O'Reilly Media, 2016, ch. 22, ``Addressing Cascading Failures.''
\bibitem{buildersretry} M. Brooker, ``Timeouts, retries, and backoff with jitter,'' \emph{Amazon Builders' Library}. [Online]. Available: \url{https://aws.amazon.com/builders-library/}
\bibitem{buildersshed} D. Yanacek, ``Using load shedding to avoid overload,'' \emph{Amazon Builders' Library}. [Online]. Available: \url{https://aws.amazon.com/builders-library/}
\bibitem{gan2019} Y. Gan et al., ``An open-source benchmark suite for microservices and their hardware-software implications for cloud and edge systems,'' in \emph{ASPLOS}, 2019.
\end{thebibliography}

\end{document}
